\documentclass[12pt,a4paper]{report}
\usepackage[T1]{fontenc}
\usepackage[margin=1in]{geometry}
\usepackage{graphicx}
\usepackage{amsmath}
\usepackage{amssymb}
\usepackage[hidelinks]{hyperref}
\usepackage{fancyhdr}
\usepackage{listings}
\usepackage{xcolor}
\usepackage{array}
\usepackage{booktabs}
\usepackage{float}
\usepackage{lmodern}
\usepackage{tikz}

% TikZ Libraries for the Verification Diagram
\usetikzlibrary{shapes.geometric, arrows.meta, positioning, calc, shapes.cylinders}

% Code listing styling
\lstset{
	basicstyle=\ttfamily\scriptsize,
	breaklines=true,
	commentstyle=\itshape\color{gray},
	keywordstyle=\bfseries\color{blue},
	stringstyle=\color{red!70!black},
	backgroundcolor=\color{lightgray!10},
	frame=single,
	captionpos=b,
	showstringspaces=false,
	numbers=left,
	numberstyle=\tiny\color{gray},
	stepnumber=1,
	numbersep=8pt,
	tabsize=4
}

% Header and Footer
\pagestyle{fancy}
\fancyhf{}
\rhead{CPEN 421 Phase 2}
\lhead{Backend Implementation}
\cfoot{\thepage}

\renewcommand{\arraystretch}{1.3}

\begin{document}

% ============================================================================
% CUSTOM TITLE PAGE
% ============================================================================
\begin{titlepage}
	\begin{center}
		\vspace*{1cm}

		\includegraphics[width=0.3\textwidth]{logo.png}\\[1cm] % Ensure you have a logo.png in your folder

		{\Large \textbf{UNIVERSITY OF GHANA}}\\[0.2cm]
		{\large School of Engineering Sciences}\\[0.2cm]
		{\large Department of Computer Engineering}\\[1cm]

		\rule{\linewidth}{0.5mm} \\[0.4cm]
		{\huge \textbf{National Emergency Response and Dispatch Platform}} \\[0.2cm]
		{\large \textbf{Phase 2: Backend Implementation}} \\[0.4cm]
		\rule{\linewidth}{0.5mm} \\[1.5cm]

		{\Large \textbf{CPEN 421: Mobile and Web Design Architecture}}\\[1cm]

		{\large \textbf{Group 5}}\\[0.5cm]

		\begin{tabular}{ll}
			\textbf{Joshua Nuku} & 11252857 \\
			\textbf{Jeffery Quansah} & 11252855 \\
		\end{tabular}

		\vfill
		
		{\large \textbf{Source Code Repository:}}\\
		\url{https://github.com/JoshNuku/emergency-dispatch-system.git}\\[0.8cm]

		{\large \textbf{Lecturer:} Dr. Gifty Osei}\\
		{\large \textbf{TA:} Nathaniel Adika}\\[0.8cm]

		\large \today

	\end{center}
\end{titlepage}

\tableofcontents
\newpage

% ============================================================================
% CHAPTER 1: IMPLEMENTATION OVERVIEW
% ============================================================================
\chapter{Implementation Overview}

\section{Technology Stack}

The backend is implemented as a distributed system utilizing a microservices architecture. To ensure high availability, type safety, and efficient concurrency, the services were built using the following core technologies:

\begin{itemize}
	\item \textbf{Language:} Go 1.22 (Golang)
	\item \textbf{Web Framework:} Gin Gonic (for high-performance HTTP REST APIs)
	\item \textbf{ORM:} GORM (for database abstraction and automated schema migrations)
	\item \textbf{Database:} PostgreSQL with the PostGIS extension for spatial queries
	\item \textbf{Message Broker:} RabbitMQ (for asynchronous, event-driven communication)
	\item \textbf{Authentication:} JWT (JSON Web Tokens) leveraging the HS256 algorithm
	\item \textbf{Security:} bcrypt for robust password hashing and Role-Based Access Control (RBAC)
\end{itemize}

\section{Codebase Structure}

The project is structured as a monorepo following the Standard Go Project Layout, ensuring a clear separation of concerns between configuration, business logic, and API routing. The complete source code can be accessed at: \url{https://github.com/JoshNuku/emergency-dispatch-system.git}

\begin{lstlisting}[language=bash, caption=Backend Monorepo Directory Structure, numbers=none]
cpen421-backend/
├── cmd/
│   ├── auth/           # Entry point for Identity Service
│   ├── incident/       # Entry point for Incident Service
│   ├── dispatch/       # Entry point for Dispatch Service
│   └── gateway/        # Entry point for Realtime WebSocket Hub
├── internal/
│   ├── models/         # GORM database entities
│   ├── middleware/     # JWT and CORS interceptors
│   ├── repository/     # Database access layer (PostgreSQL)
│   └── config/         # Environment variable loaders (.env)
├── pkg/
│   ├── rabbitmq/       # Shared AMQP publisher/consumer logic
│   └── logger/         # Centralized structured logging
├── go.mod
└── docker-compose.yml  # Infrastructure bootstrapping
\end{lstlisting}

\section{Service Breakdown}

The system features five operational services working in tandem:

\begin{enumerate}
	\item \textbf{Identity \& Authentication Service (8081):} Manages user registration, credential verification, and session security.
	\item \textbf{Emergency Incident Service (8082):} Core logic for recording emergency payloads and computing dispatch routing.
	\item \textbf{Dispatch Tracking Service (8083):} Manages the responder vehicle registry and real-time GPS telemetry updates.
	\item \textbf{Analytics \& Monitoring Service (8084):} Aggregates system metrics (e.g., response times) for operational insights.
	\item \textbf{Realtime Gateway (8085):} A stateless WebSocket hub that broadcasts RabbitMQ live events to authenticated web clients.
\end{enumerate}

\newpage

% ============================================================================
% CHAPTER 2: MICROSERVICES IMPLEMENTATION
% ============================================================================
\chapter{Microservices Implementation}

\section{Identity and Authentication Service}

This service handles the user lifecycle and enforces security boundaries. We implemented the Repository Pattern to abstract GORM operations away from the HTTP controllers, allowing for easier unit testing.

\begin{lstlisting}[language=Go, caption=User Model (GORM)]
package models

import "github.com/google/uuid"

type User struct {
	ID           uuid.UUID  `gorm:"type:uuid;primaryKey" json:"id"`
	Name         string     `gorm:"size:255;not null" json:"name"`
	Email        string     `gorm:"size:255;uniqueIndex;not null" json:"email"`
	PasswordHash string     `gorm:"size:255;not null" json:"-"`
	Role         string     `gorm:"size:50;not null" json:"role"`
	StationID    *uuid.UUID `gorm:"type:uuid" json:"station_id,omitempty"`
}
\end{lstlisting}

% Added newpage to prevent the large JWT listing from breaking awkwardly across pages
\newpage
\subsection{JWT Middleware}
Stateless authentication is strictly enforced across all protected endpoints using a custom Gin middleware. It validates the HMAC signature of incoming Bearer tokens and extracts RBAC permissions.

\begin{lstlisting}[language=Go, caption=JWT Validation Middleware]
func JWTAuth(secret string) gin.HandlerFunc {
	return func(c *gin.Context) {
		authHeader := c.GetHeader("Authorization")
		tokenStr, found := strings.CutPrefix(authHeader, "Bearer ")
		
		if !found {
			c.AbortWithStatusJSON(http.StatusUnauthorized, gin.H{"error": "Missing token"})
			return
		}

		claims := &Claims{}
		token, err := jwt.ParseWithClaims(tokenStr, claims, func(t *jwt.Token) (interface{}, error) {
			return[]byte(secret), nil
		})
		
		if err != nil || !token.Valid {
			c.AbortWithStatusJSON(http.StatusUnauthorized, gin.H{"error": "Invalid or expired token"})
			return
		}
		
		// Inject role into context for downstream controllers
		c.Set("role", claims.Role)
		c.Set("user_id", claims.UserID)
		c.Next()
	}
}
\end{lstlisting}

\newpage

\section{Emergency Incident Service}

The incident service acts as the operational core of the platform, utilizing PostGIS extensions to manage geographical data efficiently.

\subsection{PostGIS Nearest Responder Selection}
When a citizen reports an emergency, the system must autonomously identify the nearest available responder station (Police, Fire, or Ambulance). We leverage the PostGIS \texttt{ST\_Distance} and \texttt{ST\_MakePoint} functions to execute meter-based distance calculations directly within the database engine, avoiding memory-heavy application-side sorting.

\begin{lstlisting}[language=Go, caption=Spatial Query for Nearest Responder]
func (r *StationRepository) FindNearestAvailable(lat, lng float64, stationType string) (*models.ResponderStation, error) {
	var station models.ResponderStation
	
	// Base query filters for availability and type
	query := r.db.Where("type = ? AND is_available = true", stationType)
	
	// Distance calculation using PostGIS geography points
	orderClause := fmt.Sprintf(
		"ST_Distance(ST_MakePoint(longitude, latitude)::geography, ST_MakePoint(%f, %f)::geography) ASC",
		lng, lat,
	)
	
	// Execute spatial sort and fetch the top result
	err := query.Order(orderClause).First(&station).Error
	if err != nil {
		return nil, fmt.Errorf("failed to locate nearby station: %w", err)
	}
	
	return &station, nil
}
\end{lstlisting}

\newpage

% ============================================================================
% CHAPTER 3: INTER-SERVICE COMMUNICATION
% ============================================================================
\chapter{Inter-service Communication}

\section{Hybrid Communication Model}

To balance system reliability, speed, and loose coupling, the backend implements a hybrid communication architecture:

\begin{enumerate}
	\item \textbf{Synchronous (HTTP/REST):} Utilized for critical operational pathways where immediate confirmation is mandatory (e.g., the Incident Service invoking the Dispatch Service to lock a vehicle reservation).
	\item \textbf{Asynchronous (RabbitMQ):} Utilized for secondary, non-blocking operations such as telemetry updates and analytics aggregation, where eventual consistency is preferred to maximize throughput.
\end{enumerate}

\section{Message Queue Implementation}

A robust RabbitMQ publisher was implemented within the Incident service. It utilizes a \texttt{topic} exchange named \texttt{emergency.events}, allowing dynamic routing of payloads based on the event type.

\begin{lstlisting}[language=Go, caption=Event Publication Logic]
func (p *Publisher) PublishEvent(routingKey string, payload interface{}) error {
	event := map[string]interface{}{
		"event_type": routingKey,
		"payload":    payload,
		"timestamp":  time.Now().UTC().Format(time.RFC3339),
	}
	
	body, err := json.Marshal(event)
	if err != nil {
		return err
	}
	
	return p.channel.Publish(
		"emergency.events", // exchange name
		routingKey,         // routing key (e.g., incident.created)
		false, false,
		amqp.Publishing{
			ContentType:  "application/json",
			DeliveryMode: amqp.Persistent, // Ensure messages survive broker restarts
			Body:         body,
		},
	)
}
\end{lstlisting}

\section{Realtime Gateway Consumer}

To push updates to the web dashboard without polling, the Realtime Gateway maintains a continuous RabbitMQ consumption loop, fanning out JSON messages to authenticated WebSocket clients.

\begin{lstlisting}[language=Go, caption=RabbitMQ Consumer in WebSocket Gateway]
func (g *Gateway) ConsumeAndBroadcast(msgs <-chan amqp.Delivery) {
	for msg := range msgs {
		// Wrap raw RabbitMQ message into a structured UI envelope
		wsMsg := map[string]interface{}{
			"type":    msg.RoutingKey, // Evaluated by frontend reducers
			"payload": string(msg.Body),
		}
		
		// Broadcast to the thread-safe WebSocket Hub
		g.hub.BroadcastJSON(wsMsg)
		
		// Acknowledge successful processing
		msg.Ack(false)
	}
}
\end{lstlisting}

\newpage

% ============================================================================
% CHAPTER 4: DATABASE IMPLEMENTATION
% ============================================================================
\chapter{Database Implementation}

\section{Data Isolation}

Adhering strictly to microservice design principles, the architecture implements Database-per-Service isolation. This bounded context ensures that a complete failure or deadlock in the Analytics database has zero impact on the core incident reporting pipeline.

\begin{table}[H]
\centering\small
\caption{Microservice Database Specifications}
\begin{tabular}{|l|l|l|}
\hline
\textbf{Service} & \textbf{Database Schema} & \textbf{Required Extensions} \\
\hline
Authentication & \texttt{auth\_db} & \texttt{uuid-ossp} \\
\hline
Incident Management & \texttt{incident\_db} & \texttt{postgis}, \texttt{uuid-ossp} \\
\hline
Dispatch \& Telemetry & \texttt{dispatch\_db} & \texttt{postgis}, \texttt{uuid-ossp} \\
\hline
Analytics Aggregation & \texttt{analytics\_db} & \texttt{uuid-ossp} \\
\hline
\end{tabular}
\end{table}

\section{Schema Evolution \& Bootstrapping}

GORM's \texttt{AutoMigrate} feature was utilized during the application bootstrap phase to dynamically ensure that backend table structures are flawlessly synchronized with the compiled Go models. 

\begin{lstlisting}[language=Go, caption=Database Bootstrap Sequence]
func InitDB(cfg *config.Config) (*gorm.DB, error) {
	dsn := fmt.Sprintf(
		"host=%s user=%s password=%s dbname=%s port=%s sslmode=disable", 
		cfg.DBHost, cfg.DBUser, cfg.DBPass, cfg.DBName, cfg.DBPort,
	)
	
	db, err := gorm.Open(postgres.Open(dsn), &gorm.Config{
		Logger: logger.Default.LogMode(logger.Info),
	})
	if err != nil {
		return nil, err
	}
	
	// Ensure spatial capabilities are active
	db.Exec("CREATE EXTENSION IF NOT EXISTS \"postgis\"")
	db.Exec("CREATE EXTENSION IF NOT EXISTS \"uuid-ossp\"")
	
	// Synchronize schemas safely
	err = db.AutoMigrate(&models.Incident{}, &models.ResponderStation{})
	
	return db, err
}
\end{lstlisting}

\newpage

% ============================================================================
% CHAPTER 5: VERIFICATION AND TESTING
% ============================================================================
\chapter{Verification and Testing}

\section{API Verification}

To ensure endpoint robustness, all routes were extensively tested via Postman collections and automated HTTP tests. The Identity service was verified against multiple security scenarios:
\begin{itemize}
	\item \textbf{Token Issuance:} Successful registration and login mapping to valid JWT generation.
	\item \textbf{Expiration Validation:} Automatic rejection of tokens older than the 15-minute TTL.
	\item \textbf{RBAC Rejection:} Unauthorized roles (e.g., Police Admin) correctly receiving HTTP 403 Forbidden when attempting to access Hospital metrics.
\end{itemize}

\section{System Integration Test Flow}

The comprehensive lifecycle of the backend was verified through a simulated end-to-end incident dispatch flow. The architectural pathway of this successful test is illustrated in Figure \ref{fig:integration_flow}.

\begin{figure}[H]
\centering
% Added resizebox to perfectly fit the diagram inside page margins and fix overflow issues
\resizebox{\textwidth}{!}{%
\begin{tikzpicture}[
    node distance=1.5cm and 2cm,
    process/.style={rectangle, draw, fill=blue!10, text width=3cm, text centered, rounded corners, minimum height=1cm, font=\small},
    db/.style={cylinder, shape border rotate=90, draw, fill=orange!10, text width=2cm, text centered, minimum height=1.2cm, font=\scriptsize},
    queue/.style={rectangle, draw, fill=purple!10, text width=3cm, text centered, minimum height=1cm, font=\small\bfseries},
    arrow/.style={thick, ->, >=stealth}
]

    % Nodes
    \node[process, fill=green!10] (client) {1. API Client\\(Postman)};
    \node[process] (incident) [right=of client] {2. Incident Service};
    \node[db] (postgis) [below=of incident] {3. PostGIS DB\\(Spatial Query)};
    \node[queue] (rmq) [right=of incident] {4. RabbitMQ\\(Exchange)};
    \node[process] (gateway) [right=of rmq] {5. Realtime Gateway};
    \node[process, fill=teal!10] (dashboard) [below=of gateway] {6. Web Dashboard\\(WebSocket)};
    \node[process] (analytics) [above=of rmq] {Analytics Service};

    % Arrows
    \draw[arrow] (client) -- node[above, font=\scriptsize] {POST /incidents} (incident);
    \draw[arrow, <->] (incident) -- node[left, font=\scriptsize] {ST\_Distance} (postgis);
    \draw[arrow] (incident) -- node[above, font=\scriptsize] {Publish} (rmq);
    \draw[arrow] (rmq) -- node[above, font=\scriptsize] {Consume} (gateway);
    \draw[arrow] (rmq) -- node[left, font=\scriptsize] {Consume} (analytics);
    \draw[arrow] (gateway) -- node[right, font=\scriptsize] {WS Push} (dashboard);

\end{tikzpicture}%
}
\caption{Event-Driven Integration Test Pathway}
\label{fig:integration_flow}
\end{figure}

The steps executed during the test sequence were:
\begin{enumerate}
	\item \textbf{Submission:} An emergency payload (lat/lng coordinates) was posted to \texttt{/incidents}.
	\item \textbf{Processing:} The PostGIS database executed the spatial query, correctly bypassing two distant stations and returning the nearest available ambulance unit.
	\item \textbf{Communication:} An \texttt{incident.dispatched} message was reliably published to the RabbitMQ exchange.
	\item \textbf{Reception:} The Realtime Gateway consumed the queue message and pushed it to an active WebSocket connection.
	\item \textbf{Persistence:} Concurrently, the Analytics service absorbed the event to record the incident dispatch times for historical reporting.
\end{enumerate}

\section{Conclusion}

Phase 2 successfully delivered a scalable, secure, and event-driven backend architecture. The integration of GORM, PostGIS, and RabbitMQ has proven highly effective at fulfilling the platform's core requirements. The backend is fully operational and ready for Phase 3: the implementation of the web-based client interface.

\end{document}